\documentclass[12pt]{article}

\usepackage[utf8]{inputenc}
\usepackage{amsmath, amssymb, amsthm}
\usepackage{geometry}
\geometry{a4paper, margin=2.5cm}

\title{Matrice Elementare}
\author{}
\date{}

\theoremstyle{definition}
\newtheorem{definition}{Definiție}

\theoremstyle{plain}
\newtheorem{proposition}{Propoziție}

\begin{document}

\maketitle

\section*{Matrice elementare}

\begin{definition}
Fie $R$ un inel comutativ unitate și $n \ge 1$. O matrice $E \in M_n(R)$ se numește
\emph{matrice elementară} dacă se obține din matricea identitate $I_n$ printr-o singură
operație elementară pe linii (sau coloane). Există trei tipuri fundamentale:
\end{definition}

\begin{definition}[Tipul I: Permutare]
Pentru $1 \le i < j \le n$, matricea


\[
P_{ij}
\]


se obține prin permutarea liniilor $i$ și $j$ ale lui $I_n$. Este inversabilă, cu


\[
P_{ij}^{-1} = P_{ij}.
\]


\end{definition}

\begin{definition}[Tipul II: Scalare]
Pentru $a \in R^\times$ (unitate), matricea


\[
D_i(a)
\]


se obține înmulțind linia $i$ a lui $I_n$ cu $a$. Inversa este


\[
D_i(a)^{-1} = D_i(a^{-1}).
\]


\end{definition}

\begin{definition}[Tipul III: Adăugare de multiplu]
Pentru $i \ne j$ și $a \in R$, matricea


\[
T_{ij}(a)
\]


se obține adăugând $a$ ori linia $j$ la linia $i$ în $I_n$. Inversa este


\[
T_{ij}(a)^{-1} = T_{ij}(-a).
\]


\end{definition}

\begin{proposition}
Toate matricile elementare sunt inversabile. Mai mult:


\[
P_{ij}^2 = I_n, \qquad
D_i(a) D_i(b) = D_i(ab), \qquad
T_{ij}(a) T_{ij}(b) = T_{ij}(a+b).
\]


\end{proposition}

\begin{proposition}
Pentru orice $A \in M_n(R)$, aplicarea unei operații elementare pe linii este echivalentă cu
înlocuirea lui $A$ prin $EA$, unde $E$ este o matrice elementară.
Analog, operațiile pe coloane sunt de forma $AE$.
\end{proposition}

\begin{proposition}
Dacă $R$ este un PID, atunci orice matrice inversabilă $U \in GL_n(R)$ se poate scrie ca un produs finit de matrici elementare:


\[
U = E_1 E_2 \cdots E_k.
\]


\end{proposition}

\begin{proposition}
Matricele elementare generează grupul $GL_n(R)$:


\[
GL_n(R) = \langle P_{ij},\, D_i(a),\, T_{ij}(a) \rangle.
\]


\end{proposition}

\begin{proposition}[Comutații utile]
Pentru $i,j,k$ distincte și $a,b \in R$:


\[
T_{ij}(a) T_{ik}(b) = T_{ik}(b) T_{ij}(a),
\]




\[
T_{ij}(a) T_{kj}(b) = T_{kj}(b) T_{ij}(a),
\]




\[
T_{ij}(a) T_{jk}(b) = T_{jk}(b) T_{ik}(ab) T_{ij}(a).
\]


Aceste relații sunt fundamentale în reducerea matricilor la forma Smith.
\end{proposition}



\section*{Matrice elementare scrise explicit}

\subsection*{1. Matricea de permutare $P_{ij}$}

Matricea $P_{ij}$ se obține prin permutarea liniilor $i$ și $j$ ale lui $I_n$:



\[
P_{ij} =
\begin{pmatrix}
1      &        &        &        &        &        \\
       & \ddots &        &        &        &        \\
       &        & 0      & \cdots & 1      &        \\
       &        & \vdots & \ddots & \vdots &        \\
       &        & 1      & \cdots & 0      &        \\
       &        &        &        &        & \ddots \\
\end{pmatrix}
\]



unde liniile (sau coloanele) $i$ și $j$ sunt schimbate între ele.



\[
P_{ij}^{-1} = P_{ij}.
\]



\subsection*{2. Matricea de scalare $D_i(a)$}

Matricea $D_i(a)$ se obține înmulțind linia $i$ cu $a$:



\[
D_i(a) =
\begin{pmatrix}
1      &        &        &        &        \\
       & \ddots &        &        &        \\
       &        & a      &        &        \\
       &        &        & \ddots &        \\
       &        &        &        & 1
\end{pmatrix}
\]



unde $a$ apare pe poziția $(i,i)$.



\[
D_i(a)^{-1} = D_i(a^{-1}).
\]



\subsection*{3. Matricea de adăugare $T_{ij}(a)$}

Matricea $T_{ij}(a)$ se obține adăugând $a$ ori linia $j$ la linia $i$:



\[
T_{ij}(a) =
\begin{pmatrix}
1      &        &        &        &        \\
       & \ddots &        &        &        \\
       &        & 1      &        &        \\
       &        & a      & 1      &        \\
       &        &        &        & \ddots
\end{pmatrix}
\]



unde elementul $(i,j)$ este $a$, iar restul sunt ca în $I_n$.



\[
T_{ij}(a)^{-1} = T_{ij}(-a).
\]



\subsection*{4. Matricea $Q_{ij}$}

Matricea $Q_{ij}$ se obține din matricea identitate $I_n$ înlocuind blocul $2\times 2$
corespunzător liniilor și coloanelor $i$ și $j$ cu



\[
\begin{pmatrix}
0 & -1 \\
1 & 0
\end{pmatrix}.
\]



Explicit,


\[
Q_{ij} =
\begin{pmatrix}
1      &        &        &        &        &        \\
       & \ddots &        &        &        &        \\
       &        & 0      & -1     &        &        \\
       &        & 1      & 0      &        &        \\
       &        &        &        & \ddots &        \\
       &        &        &        &        & 1
\end{pmatrix}.
\]



Avem


\[
\det Q_{ij} = 1, \qquad Q_{ij}^{-1} = Q_{ij}^T =
\begin{pmatrix}
0 & 1 \\
-1 & 0
\end{pmatrix}
\text{ în blocul } (i,j).
\]







\end{document}
